\writeSectionFrame{\presentationTitle}{Das Problem}
\begin{frame}{Wir möchten sortieren}
	\begin{itemize}
 		\item BubbleSort
 		\item InsertionSort
 		\item QuickSort
 		\item SelectionSort
 		\item HeapSort
 		\item MergeSort
 		\item usw.
 	\end{itemize}
\end{frame}

\begin{frame}{Benutzer}
	\begin{center}
 		\includegraphics[scale=0.4]{\BILDERPFAD/problem1.png}
 	\end{center}
\end{frame}

\note{
 	Der Benutzer soll während der Laufzeit einen Algorithmus auswählen können. Wir brauchen
 	also eine relativ lose Kopplung zum eigentlichen Algorithmus.
}

\frameWithImageOnly{\BILDERPFAD/problem2.jpg}
\note{
    Wenn wir uns einen Onlineshop vorstellen, dann gibt es dort normalerweise verschiedene Bezahlverfahren. Die möglichen Verfahren sollen erweitert werden können, ohne dass es zu Änderungen im Programm kommt. Dabei soll das Bezahlverfahren während der Laufzeit auswählbar sein.
}

\frameWithImageOnly{\BILDERPFAD/problem3.jpg}
\note{
    Die Speicherung von Daten kann auch auf verschiedene Arten erfolgen. Wir können z.B. ein Array verwenden, eine Liste oder ein Set. Die Auswahl der Datenstruktur kann vom Benutzer abhängen oder von ganz anderen Faktoren wie z.B. der zu speichernden Datenmenge.
}